% Part: first-order-logic
% Chapter: tableaux
% Section: proving-things-quant

\documentclass[../../../include/open-logic-section]{subfiles}

\begin{document}

\olfileid[id]{fol}{tab}{prq}

\olsection{\usetoken{P}{tableau} dengan Kuantor}

\begin{ex}
Ketika menangani kuantor, kita harus memastikan bahwa syarat variabel
eigen tidak dilanggar; kadang-kadang hal ini mengharuskan kita mengubah
urutan pelaksanaan inferensi tertentu. Secara umum, sebaiknya kita
mencoba menangani terlebih dahulu aturan-aturan yang tunduk pada syarat
variabel eigen (aturan-aturan tersebut akan berada lebih atas dalam
!!{tableau} akhir).

Mari kita lihat cara membuat !!a{tableau} untuk !!{sentence}
$\lexists[x][\lnot !A(x)] \lif \lnot \lforall[x][!A(x)]$.
Seperti biasa, kita mulai dengan mencatat asumsinya:
\begin{oltableau}
[\sFmla{\False}{\lexists[x][\lnot \formula{A}(x)]\lif \lnot
    \lforall[x][\formula{A}(x)]}, just=\TAss]
\end{oltableau}
Karena !!{main operator}-nya adalah $\lif$, kita menerapkan
$\TRule{\False}{\lif}$:
\begin{oltableau}
[\sFmla{\False}{\lexists[x][\lnot \formula{A}(x)]\lif \lnot
    \lforall[x][\formula{A}(x)]}, just=\TAss, checked
  [\sFmla{\True}{\lexists[x][\lnot \formula{A}(x)]},
    just={\TRule{\False}{\lif}[1]}
    [\sFmla{\False}{\lnot\lforall[x][\formula{A}(x)]},
      just={\TRule{\False}{\lif}[1]}]
  ]
]
\end{oltableau}
Baris berikutnya yang perlu ditangani adalah baris~$2$. Kita menggunakan
$\TRule{\True}{\lexists}$. Aturan ini memerlukan !!{constant} baru;
karena belum ada !!{constant} yang muncul, kita boleh memilih sebarang
konstanta, misalnya~$a$.
\begin{oltableau}
[\sFmla{\False}{\lexists[x][\lnot \formula{A}(x)]\lif \lnot
    \lforall[x][\formula{A}(x)]}, just=\TAss, checked
  [\sFmla{\True}{\lexists[x][\lnot \formula{A}(x)]},
    just={\TRule{\False}{\lif}[1]}, checked
    [\sFmla{\False}{\lnot\lforall[x][\formula{A}(x)]},
      just={\TRule{\False}{\lif}[1]}
      [\sFmla{\True}{\lnot\formula{A}(a)}, just={\TRule{\True}{\lexists}[2]}]
    ]
  ]
]
\end{oltableau}
Sekarang kita menerapkan $\TRule{\False}{\lnot}$ pada baris~$3$:
\begin{oltableau}
[\sFmla{\False}{\lexists[x][\lnot \formula{A}(x)]\lif \lnot
    \lforall[x][\formula{A}(x)]}, just=\TAss, checked
  [\sFmla{\True}{\lexists[x][\lnot \formula{A}(x)]},
    just={\TRule{\False}{\lif}[1]}, checked
    [\sFmla{\False}{\lnot\lforall[x][\formula{A}(x)]},
      just={\TRule{\False}{\lif}[1]}, checked
      [\sFmla{\True}{\lnot\formula{A}(a)}, just={\TRule{\True}{\lexists}[2]}
        [\sFmla{\True}{\lforall[x][\formula{A}(x)]},
          just = {\TRule{\False}{\lnot}[3]}
        ]
      ]
    ]
  ]
]
\end{oltableau}
Kita memperoleh !!{tableau} tertutup dengan menerapkan
$\TRule{\True}{\lnot}$ pada baris~$4$, lalu
$\TRule{\True}{\lforall}$ pada baris~$5$.
\begin{oltableau}
[\sFmla{\False}{\lexists[x][\lnot \formula{A}(x)]\lif \lnot
    \lforall[x][\formula{A}(x)]}, just=\TAss, checked
  [\sFmla{\True}{\lexists[x][\lnot \formula{A}(x)]},
    just={\TRule{\False}{\lif}[1]}, checked
    [\sFmla{\False}{\lnot\lforall[x][\formula{A}(x)]},
      just={\TRule{\False}{\lif}[1]}, checked
      [\sFmla{\True}{\lnot\formula{A}(a)}, just={\TRule{\True}{\lexists}[2]}
        [\sFmla{\True}{\lforall[x][\formula{A}(x)]},
          just = {\TRule{\False}{\lnot}[3]}
          [\sFmla{\False}{\formula{A}(a)}, just={\TRule{\True}{\lnot}[4]}
            [\sFmla{\True}{\formula{A}(a)},
              just={\TRule{\True}{\lforall}[5]}, close
            ]
          ]
        ]
      ]
    ]
  ]
]
\end{oltableau}
\end{ex}

\begin{ex}
Mari kita lihat cara membuat !!a{tableau} untuk himpunan
\[
\sFmla{\False}{\lexists[x][!C(x,b)]},
\sFmla{\True}{\lexists[x][(!A(x) \land !B(x))]},
\sFmla{\True}{\lforall[x][(!B(x) \lif !C(x,b))]}.
\]
Seperti biasa, kita mulai dengan asumsi-asumsinya:
\begin{oltableau}
  [\sFmla{\False}{\lexists[x][\formula{C}(x,b)]}, just=\TAss
    [\sFmla{\True}{\lexists[x][(\formula{A}(x) \land \formula{B}(x))]},
      just=\TAss
      [\sFmla{\True}{\lforall[x][(\formula{B}(x) \lif \formula{C}(x,b))]},
        just=\TAss]
    ]
  ]
\end{oltableau}
Kita sebaiknya selalu menerapkan terlebih dahulu aturan yang memiliki
syarat variabel eigen; dalam kasus ini, aturan tersebut adalah
$\TRule{\True}{\lexists}$ pada baris~$2$. Karena asumsi-asumsinya memuat
!!{constant}~$b$, kita harus menggunakan konstanta lain; mari kita pilih
lagi~$a$.
\begin{oltableau}
  [\sFmla{\False}{\lexists[x][\formula{C}(x,b)]}, just=\TAss
    [\sFmla{\True}{\lexists[x][(\formula{A}(x) \land \formula{B}(x))]},
      just=\TAss, checked
      [\sFmla{\True}{\lforall[x][(\formula{B}(x) \lif \formula{C}(x,b))]},
        just=\TAss
        [\sFmla{\True}{\formula{A}(a) \land \formula{B}(a)},
          just={\TRule{\True}{\lexists}[2]}]
      ]
    ]
  ]
\end{oltableau}
Jika sekarang kita menerapkan $\TRule{\False}{\lexists}$ pada baris~$1$
atau $\TRule{\True}{\lforall}$ pada baris~$3$, kita harus menentukan
suku~$t$ yang akan menggantikan~$x$. Karena aturan-aturan ini tidak
memiliki syarat variabel eigen, kita boleh memilih sebarang suku tertutup. Dalam
beberapa kasus, kita bahkan mungkin harus menerapkan aturan itu beberapa
kali dengan~$t$ yang berbeda. Namun, sebagai pedoman umum, sebaiknya kita
memilih salah satu suku yang sudah muncul dalam !!{tableau}---dalam kasus
ini, $a$ dan~$b$---dan di sini kita dapat menduga bahwa $a$ lebih mungkin
menghasilkan cabang tertutup.
\begin{oltableau}
  [\sFmla{\False}{\lexists[x][\formula{C}(x,b)]}, just=\TAss
    [\sFmla{\True}{\lexists[x][(\formula{A}(x) \land \formula{B}(x))]},
      just=\TAss, checked
      [\sFmla{\True}{\lforall[x][(\formula{B}(x) \lif \formula{C}(x,b))]}, just=\TAss
        [\sFmla{\True}{\formula{A}(a) \land \formula{B}(a)}, just={\TRule{\True}{\lexists}[2]}
          [\sFmla{\False}{\formula{C}(a,b)}, just={\TRule{\False}{\lexists}[1]}
            [\sFmla{\True}{\formula{B}(a) \lif \formula{C}(a,b)},
              just={\TRule{\True}{\lforall}[3]}
            ]
          ]
        ]
      ]
    ]
  ]
\end{oltableau}
Kita tidak memberi tanda centang pada !!{signed formula}s di baris $1$
dan $3$, sebab kita mungkin harus menggunakannya lagi. Sekarang terapkan
$\TRule{\True}{\land}$ pada baris~$4$:
\begin{oltableau}
  [\sFmla{\False}{\lexists[x][\formula{C}(x,b)]}, just=\TAss
    [\sFmla{\True}{\lexists[x][(\formula{A}(x) \land \formula{B}(x))]},
      just=\TAss, checked
      [\sFmla{\True}{\lforall[x][(\formula{B}(x) \lif \formula{C}(x,b))]}, just=\TAss
        [\sFmla{\True}{\formula{A}(a) \land \formula{B}(a)},
          just={\TRule{\True}{\lexists}[2]}, checked
          [\sFmla{\False}{\formula{C}(a,b)}, just={\TRule{\False}{\lexists}[1]}
            [\sFmla{\True}{\formula{B}(a) \lif \formula{C}(a,b)},
              just={\TRule{\True}{\lforall}[3]}
              [\sFmla{\True}{\formula{A}(a)}, just={\TRule{\True}{\land}[4]}
                [\sFmla{\True}{\formula{B}(a)}, just={\TRule{\True}{\land}[4]}
                ]
              ]
            ]
          ]
        ]
      ]
    ]
  ]
\end{oltableau}
Jika sekarang kita menerapkan $\TRule{\True}{\lif}$ pada baris~$6$,
!!{tableau} tersebut tertutup:
\begin{oltableau}
  [\sFmla{\False}{\lexists[x][\formula{C}(x,b)]}, just=\TAss
    [\sFmla{\True}{\lexists[x][(\formula{A}(x) \land \formula{B}(x))]},
      just=\TAss, checked
      [\sFmla{\True}{\lforall[x][(\formula{B}(x) \lif \formula{C}(x,b))]},
        just=\TAss
        [\sFmla{\True}{\formula{A}(a) \land \formula{B}(a)},
          just={\TRule{\True}{\lexists}[2]}, checked
          [\sFmla{\False}{\formula{C}(a,b)},
            just={\TRule{\False}{\lexists}[1]}
            [\sFmla{\True}{\formula{B}(a) \lif \formula{C}(a,b)},
              just={\TRule{\True}{\lforall}[3]}, checked
              [\sFmla{\True}{\formula{A}(a)},
                just={\TRule{\True}{\land}[4]}
                [\sFmla{\True}{\formula{B}(a)},
                  just={\TRule{\True}{\land}[4]}
                  [\sFmla{\False}{\formula{B}(a)},
                    just={\TRule{\True}{\lif}[6]},close
                  ]
                  [\sFmla{\True}{\formula{C}(a,b)},
                    just={\TRule{\True}{\lif}[6]},close
                  ]
                ]
              ]
            ]
          ]
        ]
      ]
    ]
  ]
\end{oltableau}


\end{ex}

\begin{ex}
Kita menyusun !!a{tableau} untuk himpunan
\[
\sFmla{\True}{\lforall[x][!A(x)]}, \sFmla{\True}{\lforall[x][!A(x)]
  \lif \lexists[y][!B(y)]}, \sFmla{\True}{\lnot\lexists[y][!B(y)]}.
\]
Seperti biasa, kita menuliskan asumsi-asumsinya:
\begin{oltableau}
  [\sFmla{\True}{\lforall[x][\formula{A}(x)]}, just=\TAss
    [\sFmla{\True}{\lforall[x][\formula{A}(x)] \lif
        \lexists[y][\formula{B}(y)]}, just=\TAss
      [\sFmla{\True}{\lnot\lexists[y][\formula{B}(y)]}, just=\TAss
      ]
    ]
  ]
\end{oltableau}
Kita mulai dengan menerapkan aturan $\TRule{\True}{\lnot}$ pada
baris~$3$. Salah satu konsekuensi dari pedoman ``selalu terapkan lebih
dahulu aturan yang memiliki syarat variabel eigen'' adalah ``tunda
penerapan aturan kuantor tanpa syarat variabel eigen hingga aturan itu
diperlukan.'' Tunda pula aturan yang menghasilkan percabangan.
\begin{oltableau}
  [\sFmla{\True}{\lforall[x][\formula{A}(x)]}, just=\TAss
    [\sFmla{\True}{\lforall[x][\formula{A}(x)] \lif
        \lexists[y][\formula{B}(y)]}, just=\TAss
      [\sFmla{\True}{\lnot\lexists[y][\formula{B}(y)]}, just=\TAss, checked
        [\sFmla{\False}{\lexists[y][\formula{B}(y)]}, just={\TRule{\True}{\lnot}[3]}]
      ]
    ]
  ]
\end{oltableau}
Baris baru~$4$ memerlukan $\TRule{\False}{\lexists}$, yaitu aturan
kuantor tanpa syarat variabel eigen. Karena itu, kita menundanya dan
memilih menggunakan $\TRule{\True}{\lif}$ pada baris~$2$.
\begin{oltableau}
  [\sFmla{\True}{\lforall[x][\formula{A}(x)]}, just=\TAss
    [\sFmla{\True}{\lforall[x][\formula{A}(x)] \lif
        \lexists[y][\formula{B}(y)]}, just=\TAss, checked
      [\sFmla{\True}{\lnot\lexists[y][\formula{B}(y)]}, just=\TAss, checked
        [\sFmla{\False}{\lexists[y][\formula{B}(y)]},
          just={\TRule{\True}{\lnot}[3]},
          [\sFmla{\False}{\lforall[x][\formula{A}(x)]}, just={\TRule{\True}{\lif}[2]}]
          [\sFmla{\True}{\lexists[y][\formula{B}(y)]}, just={\TRule{\True}{\lif}[2]}]
        ]
      ]
    ]
  ]
\end{oltableau}
Kedua !!{signed formula}s baru memerlukan aturan yang memiliki syarat
variabel eigen, sehingga aturan-aturan itu harus diterapkan berikutnya:
\begin{oltableau}
  [\sFmla{\True}{\lforall[x][\formula{A}(x)]}, just=\TAss
    [\sFmla{\True}{\lforall[x][\formula{A}(x)] \lif
        \lexists[y][\formula{B}(y)]}, just=\TAss, checked
      [\sFmla{\True}{\lnot\lexists[y][\formula{B}(y)]}, just=\TAss, checked
        [\sFmla{\False}{\lexists[y][\formula{B}(y)]},
          just={\TRule{\True}{\lnot}[3]}
          [\sFmla{\False}{\lforall[x][\formula{A}(x)]}, just={\TRule{\True}{\lif}[2]},checked
            [\sFmla{\False}{\formula{A}(b)}, just={\TRule{\False}{\lforall}[5]}]
          ]
          [\sFmla{\True}{\lexists[y][\formula{B}(y)]}, just={\TRule{\True}{\lif}[2]},checked
            [\sFmla{\True}{\formula{B}(c)}, just={\TRule{\True}{\lexists}[5]}]
          ]
        ]
      ]
    ]
  ]
\end{oltableau}
Untuk menutup cabang-cabang tersebut, kita harus menggunakan
!!{signed formula}s pada baris $1$ dan~$4$. Aturan-aturan yang
bersesuaian (\TRule{\True}{\lforall} dan
\TRule{\False}{\lexists}) tidak memiliki syarat variabel eigen, sehingga
kita bebas memilih suku yang sesuai. Dalam hal ini, masing-masing suku
tersebut adalah $b$ dan~$c$.
\begin{oltableau}
  [\sFmla{\True}{\lforall[x][\formula{A}(x)]}, just=\TAss
    [\sFmla{\True}{\lforall[x][\formula{A}(x)] \lif
        \lexists[y][\formula{B}(y)]}, just=\TAss, checked
      [\sFmla{\True}{\lnot\lexists[y][\formula{B}(y)]}, just=\TAss, checked
        [\sFmla{\False}{\lexists[y][\formula{B}(y)]},
          just={\TRule{\True}{\lnot}[3]}
          [\sFmla{\False}{\lforall[x][\formula{A}(x)]},
            just={\TRule{\True}{\lif}[2]},checked
            [\sFmla{\False}{\formula{A}(b)},
              just={\TRule{\False}{\lforall}[5]}
              [\sFmla{\True}{\formula{A}(b)},
                just={\TRule{\True}{\lforall}[1]},close
              ]
            ]
          ]
          [\sFmla{\True}{\lexists[y][\formula{B}(y)]},
            just={\TRule{\True}{\lif}[2]},checked
            [\sFmla{\True}{\formula{B}(c)},
              just={\TRule{\True}{\lexists}[5]}
              [\sFmla{\False}{\formula{B}(c)},
                just={\TRule{\False}{\lexists}[4]},close
              ]
            ]
          ]
        ]
      ]
    ]
  ]
\end{oltableau}
\end{ex}

\begin{prob}
Berikan !!{tableau}s tertutup untuk hal-hal berikut:
\begin{enumerate}
\item $\sFmla{\False}{(\lforall[x][!A(x)] \land \lforall[y][!B(y)])
\lif \lforall[z][(!A(z) \land !B(z))]}$.
\item $\sFmla{\False}{(\lexists[x][!A(x)] \lor \lexists[y][!B(y)])
\lif \lexists[z][(!A(z) \lor !B(z))]}$.
\item $\sFmla{\True}{\lforall[x][(!A(x) \lif !B)]},
\sFmla{\False}{\lexists[y][!A(y)] \lif !B}$.
\item $\sFmla{\True}{\lforall[x][\lnot !A(x)]},
\sFmla{\False}{\lnot\lexists[x][!A(x)]}$.
\item $\sFmla{\False}{\lnot\lexists[x][!A(x)] \lif \lforall[x][\lnot
!A(x)]}$.
\item $\sFmla{\False}{\lnot\lexists[x][\lforall[y][((!A(x,y) \lif
\lnot !A(y,y)) \land (\lnot !A(y,y) \lif !A(x,y)))]]}$.
\end{enumerate}
\end{prob}

\begin{prob}
Berikan !!{tableau}s tertutup untuk hal-hal berikut:
\begin{enumerate}
\item $\sFmla{\False}{\lnot\lforall[x][!A(x)] \lif \lexists[x][\lnot!A(x)]}$.
\item $\sFmla{\True}{(\lforall[x][!A(x)] \lif !B)}, \sFmla{\False}{\lexists[y][(!A(y) \lif !B)]}$.
\item $\sFmla{\False}{\lexists[x][(!A(x) \lif \lforall[y][!A(y)])]}$.
\end{enumerate}
\end{prob}

\end{document}
